\section*{Conventions and scope}
\begingroup\small
We use the complex orientation of every complex manifold. The base is
$B=\PP^1_t$, with $p_1=0$, $p_2=1$, and $p_0=\infty$; $q=1/t$ at the cusp.
The regular fibre lattice is
$\Lambda=\Z\langle\widehat\gamma,\widehat u,\widehat w,\widehat\delta\rangle$.
The covector $\gamma$ is dual to $\widehat\gamma$.
In the manuscript notation $V=\Lambda^*$ and $A_j=(T_j^{-1})^t$.
The based generators have $g_1g_2g_0=1$ and finite root coordinates
$s_j\circ g_j=e^{-2\pi i/m_j}s_j$; whenever a positive geometric meridian
is used instead, we state the inversion explicitly.
The multiplicities are $m_1=3,m_2=4$; $f^*p_j=m_jS_j$ whereas $W=f^{-1}(p_0)$
is reduced. The parameter is $w=e^{-\pi i c}$; a positive $w$-loop lifts to
$c(r)=c_0-2r$, $0\leq r\leq1$.

Engel's indices $0,1,\infty$ correspond to our $1,2,0$.
His rational elliptic surface has fibres $\mathrm{IV}^*,\mathrm{III},\mathrm I_1$,
and we write $M_S=\mathcal O_S(P-O)$ for his degree-zero line bundle.
We reserve $M$ for the maximum of the period inequality correction.
The change from our homology columns to Engel's marking is
\[
P_E=\begin{pmatrix}0&-2&-1&-2\\2&-1&0&-1\\0&2&1&3\\-1&0&0&0\end{pmatrix},
\qquad \psi P_E=-\gamma,\qquad a_{j-1}=-P_Ev_j/m_j.
\]
Here $\psi$ is Engel's invariant covector. His positive root generator
is the inverse of ours; this accounts for the minus sign in the affine
translation. The conjugacy and sign checks are written in A.1.
The standard twist vectors are $v_1=(1,2,-4,0)^t$ and
$v_2=(-1,-3,3,0)^t$; general twists have
$p=12l_0-4l_1-3l_2$ and $\pi_1(X_p)=\Z/|p|$.
Our period and parameter conventions are
\[
\Pi_c=\begin{pmatrix}6\mu&\tau&1&0\\\beta^{\rm part}+c&\mu&0&1\end{pmatrix},
\quad D=\Im\beta_c-\frac{6(\Im\mu)^2}{\Im\tau},
\quad\Im c<-M=\frac3\pi\log6.
\]
The coordinate on the uniformizing upper half-plane is $z$;
$z_1,z_2$ are the finite generator fixed points, and
$M=\max_{B\setminus\{p_0\}}D_{\rm part}$.
The generator of $\Theta_7=\Z/28$ has Eells--Kuiper invariant
$1/28$, with the boundary orientation of its positive spin filling.

All homology and cohomology groups have integral coefficients unless
otherwise specified. We distinguish a literal quotient, its replacement
over a cusp neighbourhood, and an arbitrary continuous extension.
Results read only in supplied unpublished computational work are checked
in the precise form used below. We make no claim of priority.
\endgroup
